%===============================================================================
\section{Discussion}\label{sec:discussion}
%===============================================================================

\paragraph{What changes for grading practice}
The inference risk of a power-data field is not a property of the field. It depends on the background
the adversary holds, and in our data this dependence is the rule rather than the exception: the
typical field gains a third of its worst-case risk over an average background, a single background
field supplies three quarters of the amplification, and the most amplified fields---prices, renewable
forecasts, available capacity---are exactly those a market must publish. Single-field grading is
therefore not a conservative approximation of a background-aware audit; it errs in the unsafe
direction. The profile proposed here keeps the familiar form of a grade per field but reports what the
grade depends on: the single-field risk $M^{(0)}$, the amplification $\Gamma^{(K)}$, and the family of
backgrounds that attain it. In the controlled mechanism these backgrounds are the physically correct
ones---the unit output for every price, and the own-bus price before the remote one, except when the
line between them is congested. Because $\Mik$ is defined relative to the public baseline, the profile
is recomputed whenever something is released; the amortized attacker makes this cheap and
certification keeps it honest.

\paragraph{Relation to inference-graph approaches}
Propagating risk over pairwise inference relations~\cite{hu2026ignn} captures chains of
single-field inferences and physical relations known to experts, and it does not require an
attacker to be trained. It cannot represent joint inference by construction (Section~\ref{sec:res-mech}),
and its scores are not checkable by an attack experiment. Set-level inferability provides the
missing ground truth; an inference graph could be used to propose candidate backgrounds for the
scan, a combination we have not evaluated.

\paragraph{What the estimator is and is not}
The amortized attacker is not required for correctness: every decision reported here is certified
by retraining. Its role is to make the scan cheap enough to be repeated, and its accuracy matters
only through the ranking of backgrounds that certification follows. Most of its accuracy comes from solving a readout
per field set; on sets of a few fields, a random feature map already supplies most of what the readout needs, which is
why target-only masked pre-training---whose representation collapses onto the targets---adds little. Pre-training
earns its place when it has to reconstruct all columns: the representation then encodes the operating state, the
estimates improve on every metric, and the advantage grows with the size of the queried sets. A maximum over many
estimated marginals is sensitive to rare readout failures rather than to average accuracy, so safeguards against
exploding predictions matter more for $\Mik$ than for $\Vy$.

\paragraph{Limitations}
\emph{(i)} Risk is defined relative to a fixed attacker family; a stronger attacker can only raise
$\Vy$ and hence the measured combination risk, but the numerical values are family-dependent.
\emph{(ii)} The adversary is assumed to hold labelled history of the target; with little auxiliary
data, $\Mik$ as a maximum statistic is inflated by selection noise, which again argues for certified
lower bounds rather than point values. \emph{(iii)} We use random train/test splits; drift over time
may change the values of $\Mik$, although not the audit procedure. \emph{(iv)} $K=2$ nearly saturates the marginal risk on nine of ten targets but not threshold
crossings, so threshold-based statements carry their budget; parity-like targets would defeat any fixed $K$. \emph{(v)} Field roles for NEM, PJM and CAISO follow the Chinese disclosure rules rather
than local market rules; they are used to define a realistic audit, not to assess those markets.
\emph{(vi)} The mechanism case has one unit and one private parameter; it validates that the
definitions recover a known structure, not that every amplification in real data has a simple physical
reading. \emph{(vii)} The inference-graph baseline reimplements the propagation idea without the expert inputs
and attention training of~\cite{hu2026ignn}.

%===============================================================================
\section{Conclusion}\label{sec:conclusion}
%===============================================================================

We defined field-level inference risk from set-level inferability, the out-of-sample accuracy of
retrained attackers on a set of fields. The budgeted marginal risk $\Mik$ is the worst-case gain
that releasing a field brings over any background of at most $K$ fields, relative to what is
already public, and its excess over single-field risk is the background amplification $\Gamma^{(K)}$;
$\Mik$ is a threshold-free safety margin, and its critical fields are exactly the
members of small minimal unsafe sets. An amortized attacker---masked pre-training with a
closed-form readout per field set---scans all small sets in milliseconds, and retraining a few
scanned backgrounds per field certifies the result. On four power data sets with rule-based field
roles and exhaustive retraining as ground truth, and on a controlled market case with a known
mechanism, background amplification proved common, physically interpretable and invisible to
single-field and attribution scores: single-field grading recovered 13--17\% of the critical fields,
whereas the scan--certify pipeline recovered 75--88\% of them with no false positives. The profile of
single-field risk, amplification and dangerous backgrounds is the per-field quantity a
classification standard asks for when it asks how aggregation changes a field's grade.

\section*{Data and code availability}
RTS-GMLC~\cite{barrows2020rts} and the AEMO NEMWEB archive~\cite{aemo2024nemweb} are public; the PJM and CAISO series
were obtained from the operators' public data portals. Field-role tables, the code for dispatch simulation, ground-truth
retraining, the amortized attacker and all analyses will be made available upon publication.
